\begin{frame}[t]{\ev{\tagP\,\tagO}A pilot can expose shared failure classes}
\begin{center}
\begin{tikzpicture}[x=1cm,y=1cm]
\node[proposed,text width=3.6cm,minimum height=1.3cm] (a) at (0,0) {small pilot: shared parent and separate state};
\node[candidate,text width=3.6cm,minimum height=1.3cm] (b) at (4.9,0) {failure classes and their uncertainty};
\node[evidence,text width=3.6cm,minimum height=1.3cm] (c) at (9.8,0) {revised branch allocation or a larger experiment};
\draw[flow] (a)--(b);\draw[flow] (b)--(c);
\end{tikzpicture}
\end{center}
\vspace{.2cm}
{\small\textbf{Published comparison.} Nine LLM judges from seven model families were unanimous on 319 of 1,000 MNLI items. On those items, 9.1\% of the unanimous decisions were wrong.\par}
\vspace{.1cm}
{\small Agreement coexisted with error in that setting. The figure is not a fork error rate.\par}
\vspace{.1cm}
{\small A pilot can expose obvious shared failures. It cannot certify low dependence or rare false acceptance.\par}
\claim{Probe a few branches before a broad fan-out.}
\sources{Kohli, \href{https://arxiv.org/abs/2605.29800}{arXiv:2605.29800} (2026 preprint), \S5.2, ChaosNLI-MNLI. Canary requests: Dean \& Barroso, The Tail at Scale, CACM 2013.}
\note{[0:40] A small pilot comes first. It compares a shared parent with separate state and names the failure classes. Kohli reports nine judges unanimous on three hundred nineteen items, and nine point one percent of those decisions were wrong. That setting is different from forking. A pilot cannot certify low dependence. Dean and Barroso send canary requests to one or two leaves first. We close with two open problems.}
\end{frame}
